
\question
A fungus is a type of 
\begin{choices}
\correctchoice chemoheterotroph 
\choice  photoheterotroph 
\choice  photoautotroph 
\choice chemoautotroph 
\choice  autoheterotroph 
\end{choices}

\question
The transfer of fluid from the glomerulus to Bowman's capsule

\begin{choices}
\correctchoice  is mainly a consequence of blood pressure in the capillaries of the glomerulus. 
\choice is very selective as to which protein-sized molecules are transferred. 
\choice usually includes the transfer of red blood cells into Bowman's capsule. 
\choice transfers large molecules  as easily as small ones. 
\choice results from active transport. 
\end{choices}

\question
Which of the following statements is correct about diffusion in or out of a cell?

\begin{choices}
\choice It requires an expenditure of energy by the cell. 
\choice It is an active process in which molecules move from a region of lower concentration to one of higher concentration. 
\choice It requires aquaporins in the cell wall. 
\choice It requires channel proteins in the cell membrane. 
\correctchoice It is a passive process in which molecules move from a region of higher concentration to a region of lower concentration. 
\end{choices}

\question
Birds secrete uric acid as their nitrogenous waste because uric acid

\begin{choices}
\choice is readily soluble in water. 
\choice excretion allows birds to live in desert environments. 
\correctchoice requires little water for nitrogenous waste disposal, thus reducing body mass. 
\choice is metabolically less expensive to synthesize than other excretory products. 
\correctchoice is less toxic than urea. 
\end{choices}

\question
A species of rabbit that follows Bergmann's rule 

\begin{choices}
\correctchoice  would be larger in Canada than in Texas. 
\choice would be smaller in Minnesota than in Florida. 
\choice would have larger ears in Maine and smaller ears in Georgia. 
\choice would be whiter in winter and browner in summer. 
\end{choices}


\question
Which of the following is \emph{not} a major activity of the stomach?

\begin{choices}
\choice Mucus secretion. 
\choice Enzyme secretion. 
\choice Mechanical digestion. 
\correctchoice Nutrient absorption. 
\choice \ce{HCl} secretion. 
\end{choices}

\question
Which of the following is an accurate statement about thermoregulation?

\begin{choices}
\choice Endotherms are warm-blooded and ectotherms are cold-blooded.  
\choice Endothermy has a lower energy cost than ectothermy. 
\correctchoice Endotherms and ectotherms differ in their primary source of heat for thermoregulation. 
\choice Endotherms are regulators and ectotherms are conformers. 
\choice Endotherms maintain a constant body temperature and ectotherms do not. 
\end{choices}

\question
Because the foods eaten by animals are often composed largely of macromolecules, this requires the animals to have mechanisms for \rule{1in}{0.4pt} not found in other groups.

\begin{choices}
\correctchoice chemical digestion. 
\choice filtration in Bowman's capsule. 
\choice absorption of nutrients. 
\choice reabsorption of solutes from the loop of Henle. 
\choice elimination of nitrogenous waste. 
\end{choices}

\question
Which of the following contains digestive enzymes?

\begin{choices}
\correctchoice lysosome 
\choice food vacuole. 
\choice guard cells. 
\choice contractile vacuole. 
\choice central vacuole. 
\end{choices}

\question
Stomach cells are moderately well adapted to the acidity and digesting activities in the stomach by having

\begin{choices}
\choice a sufficient colony of \textit{H. pylori} to breakdown the cellulose from vegetables. 
\correctchoice a thick, mucous secretion and active replacement of epithelial cells. 
\choice a high level of bile. 
\choice secretions enter the stomach from the pancreas. 
\choice a high level of secretion of pepsin. 
\end{choices}

\newpage

\question
When the digestion and absorption of carbohydrates results in more energy-rich molecules than are immediately required by an animal, the excess is

\begin{choices}
\choice stored as starch. 
\choice eliminated in the feces. 
\correctchoice stored as glycogen or fat. 
\choice stored as cellulose or chitin. 
\choice stored in the crop. 
\end{choices}

\question
Upon activation by stomach acidity, the secretions of the gastric glands

\begin{choices}
\choice delay digestion until the food arrives in the small intestine. 
\choice include nitrogenase and amylase. 
\choice initiate the chemical digestion of lipids in the stomach.  
\choice initiate the mechanical digestion of lipids in the stomach. 
\correctchoice initiate the digestion of protein in the stomach. 
\end{choices}

\question
After ingestion by humans, the first type of macromolecules to be chemically digested by enzymes in the mouth is

\begin{choices}
\choice minerals. 
\correctchoice carbohydrates. 
\choice cholesterol and other lipids. 
\choice proteins. 
\choice nucleic acids. 
\end{choices}

\question
Which of the following membrane activities requires energy from ATP?

\begin{choices}
\choice movement of water into a cell. 
\choice movement of glucose molecules into a bacterial cell from a solution containing a higher concentration of glucose than inside the cell. 
\correctchoice movement of \ce{Na+} ions from a lower concentration in a mammalian cell to a higher concentration in the extracellular fluid. 
\choice diffusion of chloride ions across the membrane through a chloride channel. 
\choice movement of water out of a paramecium. 
\end{choices}

\question
Ingested dietary substances must cross cell membranes to be used by the body, a process known as

\begin{choices}
\choice elimination. 
\choice hydrolysis. 
\choice digestion. 
\choice ingestion. 
\correctchoice absorption. 
\end{choices}

\newpage

\question
Which of the following organs is incorrectly paired with its function?

\begin{choices}
\choice Stomach-protein digestion. 
\choice Pancreas-enzyme production. 
\correctchoice Large intestine-bile production. 
\choice Mouth-starch digestion. 
\choice Small intestine-nutrient absorption. 
\end{choices}

\question
Organisms that obtain energy by oxidizing iron compounds and obtain carbon from \ce{CO2} are called

\begin{choices}
\choice autoheterotrophs. 
\correctchoice chemoautotrophs. 
\choice photochemotrophs. 
\choice chemoheterotrophs. 
\choice photoautotrophs. 
\choice photoheterotrophs. 
\end{choices}

\question
A marine iguana is basking on a sunny rock after feeding in the cold Pacific Ocean water. It's warming its body through 

\begin{choices}
\choice convection and evaporation. 
\choice conduction and evaporation. 
\correctchoice conduction and radiation. 
\choice convection and radiation.  
\choice conduction and convection 
\end{choices}

\question
Cell walls are found in at least some representatives of  
\begin{choices}
\choice plants 
\choice protists 
\correctchoice all of the above.  
\choice bacteria 
\choice algae 
\choice fungi 
\end{choices}

\question
Small mammals like mice use more energy per unit mass than large mammals mostly because most of the energy is used for

\begin{choices}
\choice growing quickly. 
\choice reproduction. 
\choice the basal metabolic rate. 
\choice activities like hunting. 
\correctchoice thermoregulation. 
\end{choices}

\newpage

\question
Pepsin is a digestive enzyme that

\begin{choices}
\choice digests polysaccharides to monosaccharides in the small intestine. 
\choice present in bile to digest lipids in the small intestine. 
\correctchoice begins the digestion of proteins in the stomach. 
\choice secreted by the pancreas to digest proteins in the duodenum. 
\choice digests nucleic acids in the stomach. 
\end{choices}

\question
Which of the following types of molecules are the major structural components of the cell membrane?

\begin{choices}
\choice proteins and collagen. 
\choice phospholipids and cellulose. 
\choice proteins and cellulose. 
\choice nucleic acids and proteins. 
\correctchoice phospholipids and proteins. 
\end{choices}

\question
The advantage of excreting nitrogenous wastes as urea rather than as ammonia is that

\begin{choices}
\choice less nitrogen is removed from the body. 
\correctchoice urea is less toxic than ammonia. 
\choice urea can be exchanged for \ce{Na+}. 
\choice urea does not affect the osmoregulation. 
\choice urea requires more water for excretion than ammonia. 
\end{choices}

\question
Stomata will open when  

\begin{choices}
\correctchoice the guard cells have a higher solute concentration than outside the cell. 
\choice the guard cells have a lower solute concentration than outside the cell. 
\choice the pressure potential inside the guard cell is negative. 
\choice the water potential inside the guard cell is higher than outside the cell. 
\choice the solute potential inside the guard cell is positive. 
\end{choices}

\question
The osmoregulatory/excretory system of a freshwater flatworm is based on the operation of

\begin{choices}
\choice metanephridia. 
\choice Malpighian tubules. 
\choice nephrons. 
\correctchoice protonephridia. 
\choice ananephredia. 
\end{choices}
